\subsection{Targeted paired all-electron basis controls}
\label{sec:molecular-basis-controls}

A targeted seven-reaction study was selected before calculation to sample covalent, hydrogen-bonded, dispersion-bound, and torsional energy differences. Six reaction pairs are complete in the present assessment. W4-11/138 required longer SCF convergence and is not included in the paired basis comparison. The controls retain the production geometries, charges, spin multiplicities, AE Hamiltonian, D3(BJ) convention, 640/60~Ry cutoffs, and 100/434 quadrature. Both basis levels use OT with threshold $5\times10^{-7}$. Thus the intended change is the AE orbital basis, not a change in the reference reaction or integration protocol. Mean absolute errors (MAEs) remain defined on the fixed production reaction set, not this targeted selection.

Table~\ref{tab:molecular-basis-controls} reports every completed pair, including the initially obtained QZVPP ethane result. The corresponding TZVPP reactions reproduce the archived production values to within $1.4\times10^{-5}$~kcal~mol$^{-1}$. Enlarging the basis reduces the CARBHB12/4 reference error from 5.61 to 0.03~kcal~mol$^{-1}$, the G2RC/15 error from 4.98 to 2.17~kcal~mol$^{-1}$, and the WATER27/20 error from 7.16 to 2.58~kcal~mol$^{-1}$. The last comparison involves dissociation of a hydroxide--water complex and confirms substantial basis sensitivity in a charged hydrogen-bonded system. The effects on ADIM6/2 and RG18/1 are much smaller and do not uniformly improve the reference agreement. The GauXC and PySCF columns use the same reaction-resolved source data as Sec.~\ref{sec:cross-code}, retaining their own basis/core and dispersion protocols.

\begin{table}[htbp]
\centering\small
\caption{Targeted paired molecular basis controls, in kcal~mol$^{-1}$. Both bases use the same geometry, Hamiltonian, quadrature, and dispersion. The starred QZVPP ethane barrier uses the initially obtained higher-energy SCF solution. A lower-energy restart control is discussed in the text. These selected reactions do not define a revised 70-reaction MAE.}
\label{tab:molecular-basis-controls}
\begin{tabular}{lrrrrr}
\toprule
Reaction & Reference & TZVPP & QZVPP & GauXC AE & PySCF\\
\midrule
CARBHB12/4 & 9.967 & 15.577 & 9.939 & 10.388 & 10.369\\
ADIM6/2 & 1.990 & 2.088 & 2.321 & 1.852 & 1.860\\
BHROT27/1 & 2.730 & 2.836 & 28.009$^{*}$ & 2.889 & 2.892\\
RG18/1 & 0.080 & 0.134 & 0.130 & 0.136 & 0.133\\
G2RC/15 & -39.430 & -44.408 & -41.605 & -36.070 & -35.925\\
WATER27/20 & 26.600 & 33.764 & 29.183 & 26.309 & 26.250\\
\bottomrule
\end{tabular}
\end{table}


The initial QZVPP BHROT27/1 barrier of 28.009~kcal~mol$^{-1}$ differs strongly from both the 2.730 reference and the 2.836 TZVPP result, although both conformers satisfy the SCF threshold, normal termination, and electron-count checks. Restarting the eclipsed conformer from the converged staggered-conformer orbitals at unchanged geometry, basis, and quadrature reaches a lower-energy solution and gives a barrier of 2.797~kcal~mol$^{-1}$. By contrast, increasing the radial/Lebedev quadrature to 200/974 with atomic initial guesses gives 27.553~kcal~mol$^{-1}$. Thus tighter quadrature alone does not resolve the discrepancy, and the large initial barrier is sensitive to electronic-solution selection rather than being an intrinsic consequence of enlarging the basis. The initial results and the separate state-selection controls are retained for comparison. No aggregate accuracy is inferred from this deliberately selected basis study, and no favourable QZVPP replacements are mixed into the uniform 70-reaction dataset.
